\begin{proposition}[5.12.8]
\label{IV.2.s5.12.8-fr}
(lema de Hironaka). Sea $A$ un anillo local noetheriano reducido,
equidimensional y catenario. Supongamos además que, para todo ideal primo
minimal $\mathfrak q_i$ de $A$, el anillo $B_i=A/\mathfrak q_i$ sea tal que
$B_i^{(1)}$ \textup{(notación de (5.10.17))} sea una $B_i$-álgebra finita \textup{(lo que
ocurrirá, por ejemplo, cuando $A$ posea una de las dos propiedades (i), (ii) de
(5.11.6))}. Sea $t$ un elemento de $A$ que no sea divisor de cero en $A$, y tal
que:
\begin{enumerate}
\item[\textup{(i)}] El $A$-módulo $A/tA$ solo tenga un ideal primo asociado no inmerso
$\mathfrak p$ \textup{(necesariamente de altura 1 en virtud del Hauptidealsatz (\textbf{0},
16.3.2), pero no se supone que sea el único ideal primo asociado a $A/tA$)}.
\item[\textup{(ii)}] El ideal $\mathfrak pA_{\mathfrak p}$ de $A_{\mathfrak p}$ esté
engendrado por $t/1$.
\item[\textup{(iii)}] El anillo $A/\mathfrak p$ sea integralmente cerrado.
\end{enumerate}
En estas condiciones, el anillo $A$ es íntegro e integralmente cerrado, y se
tiene $\mathfrak p=tA$.
\end{proposition}

Si $Z=V(tA)$, la hipótesis (i) implica que $Z$ es una parte cerrada
irreducible de $X=\operatorname{Spec}(A)$, cuyo punto genérico $z$ es tal que
$\mathfrak j_z=\mathfrak p$. La hipótesis (ii) muestra que el ideal maximal
$\mathfrak pA_{\mathfrak p}$ del anillo local noetheriano $A_{\mathfrak p}$
está engendrado por un solo elemento, por lo que $A_{\mathfrak p}$ es un
\emph{anillo de valoración discreta} para el cual $t/1$ es un uniformizante (Bourbaki,
\emph{Alg. comm.}, chap. VI, § 3, n\textsuperscript{o} 6, prop. 9);
$A_{\mathfrak p}/tA_{\mathfrak p}$ es el cuerpo residual de
$A_{\mathfrak p}$ y, por tanto, un $A_{\mathfrak p}$-módulo de longitud 1. En
virtud de (3.4.2), esto implica que $Z$ está contenido en \emph{una sola} de las
componentes irreducibles $X_i$ de $X$; para cualquier otra componente
irreducible $X_j$ de $X$, se tiene, por consiguiente,
$\dim(X_j\cap Z)<\dim(Z)$. Por otra parte, como $t$ no es divisor de cero en
$A$, no está contenido en ninguno de los $\mathfrak q_i$ y se tiene, por
consiguiente, (5.1.8)
$\operatorname{codim}(X_i\cap Z,X_i)
=\operatorname{codim}(X_j\cap Z,X_j)=1$,
cualquiera que sea $j\neq i$. Como se supone que $A$ es biequidimensional, la
relación $\dim(X_i\cap Z)\neq\dim(X_j\cap Z)$ implicaría entonces
$\dim(X_i)\neq\dim(X_j)$ (\textbf 0, 14.3.3.1), lo cual es absurdo. Se ve, por tanto,
que solo hay \emph{un ideal primo minimal} en $A$; como $A$ es reducido por hipótesis,
esto implica que es \emph{íntegro}. Observemos ahora que, puesto que $A^{(1)}$ es una
$A$-álgebra finita, es un anillo noetheriano semi-local, y todo ideal primo
asociado no inmerso de $A^{(1)}/tA^{(1)}$ es, por tanto, de altura 1 en virtud
del Hauptidealsatz (\textbf 0, 16.3.2); ahora bien, para tal ideal $\mathfrak r$,
$\mathfrak r\cap A$ es de altura 1 en virtud de (5.10.17, (iv)) y, como
contiene a $tA$, no puede ser sino $\mathfrak p$ en virtud de la hipótesis (i).
Por otra parte, $A^{(1)}$ está contenido en la clausura integral $A'$ de $A$;
si se pone $S=A\setmin\mathfrak p$, la clausura integral de $A_{\mathfrak p}$ es
$S^{-1}A'$ (Bourbaki, \emph{Alg. comm.}, chap. V, § 1,
n\textsuperscript{o} 5, prop. 17), por lo que $S^{-1}A'=A_{\mathfrak p}$, ya
que $A_{\mathfrak p}$ es un anillo de valoración discreta; de ello se deduce
(\emph{loc. cit.}, § 2, n\textsuperscript{o} 1, lemme 1) que existe \emph{un solo}
ideal primo $\mathfrak p'$ de $A'$ por encima de $\mathfrak p$ y, \emph{a fortiori},
un solo ideal primo $\mathfrak p^{(1)}$ de $A^{(1)}$ por encima de
$\mathfrak p$, y se tiene
$A_{\mathfrak p}=A'_{\mathfrak p'}
=A^{(1)}_{\mathfrak p^{(1)}}$.
Observemos ahora que $A^{(1)}$ verifica $(S_2)$ (5.10.17, (i)) y, como $t$ no
es divisor de cero en $A^{(1)}$, $A^{(1)}/tA^{(1)}$ no tiene ideales primos
asociados

\oldpage[IV-2]{130}

inmersos (5.7.7). El ideal $\mathfrak p^{(1)}$ es, por tanto, el único ideal
primo asociado a $A^{(1)}/tA^{(1)}$; dicho de otro modo, $tA^{(1)}$ es un ideal
$\mathfrak p^{(1)}$-primario de $A^{(1)}$; esto equivale a decir que
$tA^{(1)}$ es la imagen inversa en $A^{(1)}$ de
$tA^{(1)}_{\mathfrak p^{(1)}}=tA_{\mathfrak p}$, y este último es el ideal
maximal de $A_{\mathfrak p}$ según (ii); por tanto,
$tA^{(1)}=\mathfrak p^{(1)}$ es primo en $A^{(1)}$. Por otra parte,
$A^{(1)}/\mathfrak p^{(1)}$ es finito sobre $A/\mathfrak p$ y tiene el mismo
cuerpo de fracciones (a saber, el cuerpo residual de
$A^{(1)}_{\mathfrak p^{(1)}}=A_{\mathfrak p}$), por lo que es idéntico a
$A/\mathfrak p$ en virtud de la hipótesis (iii). Se puede escribir entonces
$A^{(1)}=A+\mathfrak p^{(1)}=A+tA^{(1)}$ y, como $t$ está contenido en el
radical del anillo noetheriano semi-local $A^{(1)}$, el lema de Nakayama
implica que $A^{(1)}=A$, de donde $\mathfrak p^{(1)}=\mathfrak p=tA$. Pero,
como $A$ es catenario y $A/tA$ es íntegro e integralmente cerrado en virtud de
la hipótesis (iii), se concluye de (5.12.7) que $A$ es integralmente cerrado.
C.Q.F.D.

\begin{remark}[5.12.9]
\label{IV.2.s5.12.9-fr}
La conclusión de (5.12.8) deja de ser válida si ya no se supone que $A$ sea
equidimensional. Por ejemplo, sean $k$ un cuerpo, $B$ el anillo de polinomios
$k[X,Y,Z]$ en tres indeterminadas, $C$ el anillo
$B/\mathfrak r_1\mathfrak r_2$, donde $\mathfrak r_1=BZ$ y
$\mathfrak r_2=BX+BY$; sean $\mathfrak m$ el ideal maximal $BX+BY+BZ$,
$\mathfrak n=\mathfrak m/\mathfrak r_1\mathfrak r_2$ su imagen en $C$, y $A$
el anillo local $C_{\mathfrak n}$. Por último, sea $t_0$ la imagen en $C$ del
elemento $Y+Z$ de $B$, y $t$ su imagen en $A$
($\operatorname{Spec}(C)$ está, por tanto, formado por un plano $P$ y una
recta $D$ no contenida en este plano, mientras que
$\operatorname{Spec}(C/t_0C)$ está formado por una recta $D'$ contenida en $P$
y que pasa por el punto $D\cap P$). El anillo $A$
es reducido y catenario (por ser cociente de un anillo regular), y sus ideales
primos minimales $\mathfrak q_1$, $\mathfrak q_2$ son las imágenes de
$\mathfrak r_1$, $\mathfrak r_2$. El $A$-módulo $A/tA$ solo tiene un ideal
primo asociado no inmerso $\mathfrak p$, imagen de $BY+BZ$,
$\mathfrak pA_{\mathfrak p}$ está engendrado por $t/1$ y $A/\mathfrak p$ es
isomorfo a $k[X]$; pero $A$ no es íntegro.
\end{remark}
